\section{Statement of the problem}

Considering a plant of an affine switched system described by: 

\noindent
\begin{equation}
    \label{eq_continuous_time_dynamics}
    \dot{x}(t) = A_{\sigma(t)} x(t) + b_{\sigma(t)}
\end{equation}

\begin{figure}[h]
    \centering
    \includegraphics[]{Cap4/fig/desenho_cap4_3.eps}
    \caption{Switching signals}\label{cap2_linearizacao}
\end{figure}

We assume the existence of a periodic switching signal $\bar{\sigma}(t)$
associated with a periodic trajectory $\bar{x}(t)$, as obtained in 
Chapter \ref{cap2-determination-of-target-periodic-traj}. Our objective is to introduce perturbations 
in the switching instants such that the trajectory $x(t)$ converges to the reference trajectory $\bar{x}(t)$.

Figure \ref{cap2_linearizacao} displays the $j$th cycle of the switching signals $\bar{\sigma}(t)$ and $\sigma(t)$,
which comprise $N$ switching actions. Within this cycle, the $i$th switching times of $\bar{\sigma}(t)$ and 
$\sigma(t)$ are denoted by $\bar{t}_{j, i}$ and $t_{j, i}$, respectively.
The relation between $\bar{t}_{j, i}$ and $t_{j, i}$ is given by

\noindent
\begin{equation}
    t_{j, i} = \bar{t}_{j, i} + \delta t_{j, i}
\end{equation}

To do this using predictive control, we need to find a linear equation that relates the switching times and the resulting system at the end of one cycle.